% Appendix A: Changes from Earlier Versions. Printed after the references when \ifappendix is true (main.tex). The
% build main_noappendix.tex leaves it out. Version 4 adds a subsection on its corrections to version 3.
\section{Changes from Earlier Versions}\label{app:changes}
% arXiv versions 1 (28 September 2026) and 2 (29 September 2026) have the same source files and the same PDF text,
% apart from the arXiv stamp. Version 2 changed only the abstract of the arXiv listing, where "this http URL"
% (jax.vmap, jax.pmap) became "JAX's vmap" and "JAX's pmap".
This paper is version 4 of arXiv:2609.34818. Version 2 changed only how the abstract on the arXiv page wrote two JAX
function names. The paper itself is the same in versions 1 and 2, so every claim listed here appears in both. % src: arXiv listings of 2609.34818 as read on 2026-11-23 ("[v1] Mon, 28 Sep 2026", "[v2] Tue, 29 Sep 2026"), same e-print source files in v1 and v2
% Independent check of part 3 (A3): the title of version 1, as in paper/preprint/source/main.tex (pdftitle).
Version 1 had the title \vonequote{Accelerator Choice Is Not Enough: AlphaFold2 Inference on Cloud TPUs.}
% Page counts from pdfinfo: the PDF of arXiv:2609.34818v1 has 22 pages, and paper-v3/Pazienza_ElBani_AlphaFold2_TPU_v3.pdf
% at the tag arxiv-v3 has 10 pages.
Version 3 rewrote the paper. It introduced the current title, shortened the paper from 22 pages to 10, reorganized the sections, and replaced the text of version 1 with a new text based on a re-analysis of the committed result files. % src: PDF of arXiv:2609.34818v1 ("Pages: 22"); PDF of version 3 at the tag arxiv-v3 ("Pages: 10")
Version 4 corrects version 3 after a further review. We ran no new experiment for versions 3 and 4. We checked the claims and numbers of version 1 against the committed result files. This appendix
lists the claims that changed and the reason for each.

\subsection*{Corrections in Version 4}
The note of Table~\ref{tab:experiments} now lists the runs of the baseline configuration on the TPU. Version 3 wrote that
this configuration ran once in each session.
At 1000 residues, version 3 gave the speedup on eight chips as 7.91 times the throughput of one chip and the parallel efficiency as 98.9\%. % src: recomputed: SC.speedup_8over1_L1000, SC.eff_8_L1000_pct; results/sweep/scaling_law_data.json -> grid[3].length
The one-chip throughput at this length has two digits, so this version gives ranges (Section~\ref{sec:res-pmap} and Table~\ref{tab:grid}).
Version 3 printed the coefficients of the power-law fit with two digits. These coefficients do not reproduce the
deviations that version 3 gave, so this version prints four digits.
Version 3 gave the minimum ratio between the CPU runtime and one TPU chip as more than 450. % src: recomputed: V2.min_tpu_over_cpu_with_max_tpu
With the rounding of the TPU time, the bound that holds is 449 (Section~\ref{sec:res-platforms}). % src: recomputed: V4.min_tpu_over_cpu_worst
Version 3 claimed that the copy of the parameters could not explain all of the extra time per call.
The extra times, however, are differences of single runs, and the script also copies the inputs inside the timed region (Section~\ref{sec:res-pmap}).
Version 3 did not say that the logits averaged by the external ensemble are zero with
random weights (Section~\ref{sec:res-auto}).
Version 3 gave the throughput of \texttt{vmap} with $B$ = 4 as 0.73 times the throughput of $B$ = 1. % src: recomputed: V4B.vmap_B4_rel_B1_v3; results/sweep/batching_sweep.json -> [*].batch_size
This ratio came from the rounded throughputs. From the recorded times, the ratio is 0.72 (Section~\ref{sec:res-default}). % src: recomputed: V4B.vmap_B4_rel_B1

\subsection*{Ratios Between Platforms (Section~\ref{sec:res-platforms})}
Version 1 reported a ratio of 451 between one TPU chip and the CPU runtime. That version also gave a ratio of 27.8 between one TPU chip and the T4. % src: recomputed: B12.tpu_over_cpu, B11.tpu_over_gpu
Its limitations section said that the Colab values came from \vonequote{the same measurement campaign as the TPU
numbers.}\ Our records do not support this claim.
The T4 and the CPU runtime ran in two separate Colab sessions on August 8, 2026, at 04:50 UTC and 05:19 UTC, respectively. % src: recomputed: AUG.gpu.first_epoch_utc, AUG.cpu.first_epoch_utc (A-02, H-02)
The baseline run on the TPU did not record its date. The result of this run was committed on August 8, 2026, at 17:23 UTC. % src: recomputed: V2.first_commit
The three values also came from different copies of the script (Table~\ref{tab:platforms}).
In the session of September 15, 2026, at 19:01 UTC, the T4 took a mean of 6.57~s per steady-state call, against 13.086~s in the session of August 8, 2026, at 04:50 UTC. % src: recomputed: SEP.gpu.0915.mean, B.gpu.steady, V4C.sep_gpu_0915_start, AUG.gpu.first_epoch_utc
In the three CPU sessions of September 15 to 17, 2026 (Table~\ref{tab:platforms}), the CPU runtime took means of 348.9~s to 359.3~s, against 212.113~s in the session of August 8, 2026, at 05:19 UTC. % src: recomputed: V2.sep_cpu_steady_min, V2.sep_cpu_steady_max, B.cpu.steady, V2.sep_dates, AUG.cpu.first_epoch_utc
This version gives only the smallest ratio over all sessions (Section~\ref{sec:res-platforms}).
In every session, the mean steady-state call on the T4 took more than 13.9 times as long as the longest recorded steady-state call on one TPU chip (0.471~s). % src: recomputed: V2.min_tpu_over_t4_with_max_tpu, V4.min_tpu_over_t4_worst, TPU.single_query_steady_min_max
On the CPU runtime, the mean steady-state call took more than 449 times as long. % src: recomputed: V4.min_tpu_over_cpu_worst
Version 1 also gave a ratio between the T4 and the CPU runtime. We removed this ratio.

\subsection*{Speedup on Eight Chips (Section~\ref{sec:res-pmap})}
Version 1 reported that multi-query \texttt{pmap} on eight chips reached 6.92 times the throughput of the baseline run on the TPU, with a parallel efficiency of 86.5\%. % src: recomputed: MQ.speedup_file, MQ.efficiency_pct
The two runs differed in input and in timed region. This version takes the speedup only from the matched grid. There,
every point used the same script, input family, and timed region.
Eight chips gave 6.53 to about 7.9 times the throughput of one chip. % src: recomputed: V2.speedup_8over1_min, V2.speedup_8over1_max
Version 1 explained the lower parallel efficiency at 100 residues by a fixed overhead per call. This version gives the extra time per call and leaves its cause open. % src: results/sweep/scaling_law_data.json -> grid[0].length
Version 1 also wrote that \vonequote{distinct per-chip HBM allocations confirm the work is distributed.}\ Memory shows
only that every chip held data. The division of the work follows from \texttt{jax.pmap}, which runs one replica per
device~\cite{jaxpmap}.

\subsection*{First Call and Profiler Teardown (Section~\ref{sec:res-first})}
In every August version that we have, the single-query script runs the first call with the profiler.
The version of commit 86fca03 (September 15, 2026, 17:57 UTC) adds an option that turns the profiler off. % src: recomputed: V4C.commit_86fca03
 The TPU runs did not record the duration of the profiler teardown.
In the CPU sessions of September 15, 2026, at 18:04 UTC, September 16, 2026, at 12:37 UTC, and September 17, 2026, at 08:52 UTC, the profiler teardown took 31.7~s to 31.9~s. % src: recomputed: V2.sep_cpu_tail_min, V2.sep_cpu_tail_max, SEP.cpu.0915.profiler_tail_s, V4C.sep_cpu_0915_start, V4C.sep_cpu_0916_start, V4C.sep_cpu_0917_start (M-21, P-03)
In the T4 session of September 15, 2026, at 19:01 UTC, the profiler teardown took 23.7~s. % src: recomputed: SEP.gpu.0915.profiler_tail_s, V4C.sep_gpu_0915_start
Version 1 reported that the TPU first call took 59.1 times as long as the steady-state call. That version described this ratio as uncorrected. % src: recomputed: B15.tpu_first_over_steady (H-09, D-02)
This version reports the first call as 27.78~s and gives no ratio. % src: recomputed: B.tpu.first
Table~\ref{tab:platforms} marks that this call probably includes the profiler teardown.
Version 1 also read the compilation-cache experiment as showing that more than half of the first call survives the
cache. Both first calls of that experiment include the profiler teardown, so this version does not attribute the 15.19~s of the second process to compilation. % src: results/sweep/compilation_cache.json -> warm.first_predict_seconds (P-05, P-07)
For the T4 and for the CPU runtime, version 1 subtracted the profiler teardown from the first calls of August 8, 2026. % src: recomputed: AUG.gpu.first_epoch_utc
This version uses the runs of September 15 to 17, 2026, without the profiler instead. % src: recomputed: V2.sep_dates

\subsection*{Removed Profiler Trace}
Version 1 reported a trace of one TPU first call.
In the profile of that call, a JAX function that traces and compiles had 12.55~s of self time out of 16.56~s. This share is about 76\%. % src: profiling/trace_analysis.md ("16.56s", "12.55s"); recomputed: TR.share_pct
We did not keep the trace file, and only a written analysis of it remains. No reader can check the value, so we
removed it.

\subsection*{Automatic Partitioning (Section~\ref{sec:res-auto})}
Version 1 named the automatic-partitioning experiment \vonequote{automatic sharding} and discussed it in terms of GSPMD.
It cited the GSPMD paper for the claim that a partitioner without annotations falls back to replication. That paper
does not make this claim. The runs did not record which partitioner ran. This version reports the experiment as a run
without sharding annotations. It does not say whether the call ran on one chip or on every chip.

\subsection*{Cost Comparison (Section~\ref{sec:res-cost})}
Version 1 found that the slice with one chip working cost about as much per query as the T4.
The values were \$1.25 and \$1.27 per 1000 queries. % src: recomputed: COST.pod.per_1k, COST.gpu.per_1k ("per 1000 queries") (A-04, I-06)
This parity held only for the T4 time of the session of August 8, 2026. % src: recomputed: AUG.gpu.first_epoch_utc
The T4 price came from a survey that version 1 did not cite.
The CPU price was our estimate for an instance comparable to the Google Cloud machine type n2-standard-4, with four virtual CPUs. % src: PDF of arXiv:2609.34818v1 ("approximately $0.19/hour for an instance comparable to n2-standard-4, is our own estimate")
In the sessions of September 15 to 17, 2026, the CPU runtime had two virtual CPUs, and the sessions of August 8, 2026, did not record this number. % src: recomputed: V2.sep_dates, SEP.cpu.0915.logical_cpus, AUG.cpu.first_epoch_utc
Version 1 also priced hypothetical slices of one, two, and four chips in a table. This version prices only the TPU, at Google's list rate of \$1.20 per chip-hour. % src: recomputed: V2.price_chip_hour (C-04, C-11, C-12)
It gives one conditional example for a chip rented alone.

\subsection*{Removed \AFthree{} Results}
Version 1 had a section and an appendix on \AFthree{}. The release we tested did not accept the TPU as a platform.
The \AFthree{} runs do not address the question of this paper. That question concerns the TPU slice. We removed the
section, the appendix, and all the \AFthree{} numbers.

\subsection*{Other Corrections}
Version 1 stated that the power-law fit used a sequence-length sweep from 100 to 1000 residues. The fit in fact uses the matched grid. % src: results/sweep/scaling_law_data.json -> grid[*].length
The sequence-length sweep covered 60 to 500 residues. % src: results/sweep/scaling_law_data.json -> grid[*].length; results/sweep/sequence_length_sweep.json -> [*].num_residues (L-01)
Version 1 wrote that none of the cited work reported TPU experiments, but that version cited a model trained on TPUs.
Version 1 described a controlled comparison with identical code on the three platforms.
AlphaFold2's source ran unmodified, but the benchmark scripts differed between the platforms and between the sessions (Table~\ref{tab:platforms}).
Version 1 called the copy of the script in the notebooks of August 8, 2026, uncommitted. % src: recomputed: AUG.gpu.first_epoch_utc

The notebooks in commit 2755250 hold this copy. % src: recomputed: V2.first_commit (M-30)
Seven works cited in version 1 do not make the statements attributed to them in that version. We corrected these
statements or removed the citations. Version 1 treated the files in \texttt{results/sweep/} as raw results. These files
are summaries of the Job output (Section~\ref{sec:records}). The abstract now states that the weights were random. We
also renamed three terms. A \vonequote{pod} is now a slice, a \vonequote{protein} is a query, and \vonequote{automatic
sharding} is automatic partitioning.
